% =====================================================================
\chapter{Part B: System Model and Proposed Methods}
\label{ch:partb-methods}
\thispagestyle{plain}
% =====================================================================

\section{Introduction}

Part~A reduces the number of bits each cell puts on the fronthaul. Part~B deals with what happens when even the compressed load is larger than a shared link, and decides how that link is divided among the cells. We built this work in stages, and each stage was a response to a weakness found in the one before it (Figure~\ref{fig:b-roadmap}):
\begin{itemize}
  \item \textbf{B1: the KKT newsvendor rule.} We forecast each cell's \emph{deadline-feasible rate} as a range and split the link so that every cell carries the same priority-weighted risk. This rule is our NCC 2027 submission. Among the 12 reference controllers of that comparison it was the best on two traffic families. That set did not yet include the strongest baselines of the next step, and those baselines do beat it.
  \item \textbf{Extended study.} Here we wrote the problem as a POMDP, developed a Markov-fluid theory of the deadline-feasible rate and added a third traffic family. We also implemented the two strongest baselines we could find in queueing control and policy search, namely predictive MaxWeight and a policy learned by the cross-entropy method.
  \item \textbf{B2: a Bayesian, link-aware Transformer forecaster (LCQT).} This forecaster is designed for the shared link. It has a Bayesian mixture layer and a differentiable KKT layer.
  \item \textbf{B3: the back-pressure newsvendor (BPN).} This is a two-stage policy that contains both strongest baselines exactly. It orders cells by back-pressure and then adds newsvendor margins, and it beats both baselines in all 10 scenarios of three traffic families.
\end{itemize}

\begin{figure}[H]
\centering
\begin{tikzpicture}[font=\footnotesize,>=Stealth,
   s/.style={draw,rounded corners,align=center,text width=3.15cm,minimum height=2.2cm,inner sep=4pt}]
  \node[s,fill=iitblue!10] (b1) {\textbf{B1}\\KKT newsvendor rule\\[2pt]\scriptsize $\rstar$ target, GBM quantiles, closed-form split};
  \node[s,fill=orange!12,right=0.45cm of b1] (ext) {\textbf{Extended study}\\theory + baselines\\[2pt]\scriptsize POMDP, Markov-fluid bounds, Family M, pred.\ MaxWeight, CEM};
  \node[s,fill=purple!10,right=0.45cm of ext] (b2) {\textbf{B2}\\Bayesian LCQT\\[2pt]\scriptsize patch + cross-cell Transformer, mixture, diff.\ KKT};
  \node[s,fill=good!15,right=0.45cm of b2,draw=good!60!black,thick] (b3) {\textbf{B3}\\Back-pressure newsvendor\\[2pt]\scriptsize contains both strongest baselines; wins 10/10};
  \draw[->,thick] (b1) -- (ext);
  \draw[->,thick] (ext) -- (b2);
  \draw[->,thick] (b2) -- (b3);
\end{tikzpicture}
\caption{Development of the Part~B controllers. Each step keeps the structure that works and adds what the previous step was missing.}
\label{fig:b-roadmap}
\end{figure}

\section{Problem Setting and System Model}
\label{sec:b-model}

\subsection{Cells, link and decision}

Eight 7-2x cells share one 25~Gbit Ethernet aggregation link (Figure~\ref{fig:shared}). Three of them are low-latency (LL) cells with a deadline of 2~ms and priority weight 3. The other five are eMBB cells with a deadline of 10~ms and weight 1. We reserve 3\,\% of the link for the control, management and synchronisation planes, which leaves a capacity of $C = 12.125$~Mbit per 0.5~ms slot. The peak demands of the cells add up to $2.44\,C$. This over-subscription is deliberate. Operators do the same in practice to gain from statistical multiplexing.

A coordinator acts every $T = 20$ slots (10~ms). At each decision it sets a budget $r_i$ (bits per slot) for every cell, and that budget stays fixed for the whole interval. The coordinator sees the cells' queues, head-of-line ages, demand history and spectral efficiency, but only after a telemetry delay of $\tau = 2$ slots (1~ms). Each budget must therefore cover a lead time of $H = T + \tau$ slots that the coordinator has not observed (Figure~\ref{fig:timing}).

\begin{figure}[H]
\centering
\begin{tikzpicture}[font=\footnotesize,>=Stealth]
  \draw[->,thick] (0,0) -- (13,0) node[right]{slots};
  \foreach \x in {0,0.5,...,12.5} \draw (\x,-0.07) -- (\x,0.07);
  \draw[fill=orange!25] (2.0,0.25) rectangle (3.0,0.75); \node at (2.5,0.5) {$\tau$};
  \draw[fill=iitblue!15] (3.0,0.25) rectangle (13.0,0.75); \node at (8.0,0.5) {budgets $r_i$ held for $T = 20$ slots (10\,ms)};
  \draw[->,thick,accent] (2.0,1.4) -- (2.0,0.8); \node[text=accent,align=center] at (2.0,1.75) {telemetry\\snapshot};
  \draw[->,thick,good!60!black] (3.0,1.4) -- (3.0,0.8); \node[text=good!50!black,align=center] at (3.6,1.75) {decision\\epoch $t_k$};
  \draw[decorate,decoration={brace,amplitude=5pt,mirror}] (2.0,-0.3) -- (13.0,-0.3) node[midway,below=6pt]{lead time $H = T + \tau$: arrivals unseen when deciding};
\end{tikzpicture}
\caption{Timing of one control interval. The budgets are fixed on 1~ms-old telemetry and must cover the next $H = T+\tau$ slots of unseen traffic.}
\label{fig:timing}
\end{figure}

We did not choose to decide in advance. Real deployments force it. Near-real-time RIC and E2 control loops run at 10~ms or slower, and traffic shaping in a switch is not reconfigured every slot. The DU schedulers of different cells also have no way to consult each other inside the sub-millisecond processing pipeline.

\subsection{Fronthaul demand}

With $b$-bit BFP, one PRB-layer costs the following number of fronthaul bits per slot:
\begin{equation}
  \beta(b) = (12\cdot 2\cdot b + 8)\cdot 14\cdot\kappa,
  \label{eq:beta}
\end{equation}
where $\kappa = 1.08$ is the packet overhead. This gives $\beta(9) = 3387$ bits, and 6-, 12- and 14-bit mantissas scale this by $0.68$, $1.32$ and $1.54$. \emph{Part~A enters Part~B at this point}, because the compression setting fixes $\beta$ and with it the load. A cell with user bits $a_i(t)$ and spectral efficiency $\eta_i(t)$ (bits per PRB-layer) then has fronthaul demand
\begin{equation}
  d_i(t) = \lceil a_i(t)/\eta_i(t)\rceil\,\beta,
\end{equation}
so the fronthaul load of two cells with the same user traffic can differ by a factor of 2--3. A full cell (273 PRBs, 4 layers) peaks at about 7.4~Gbit/s.

\subsection{Constraints, queues and metric}

The budgets satisfy
\begin{equation}
  \sum_i r_i \le C, \qquad 0 \le r_i \le \mathrm{cap}_i = 273\,L_i\,\beta .
\end{equation}
Each cell keeps a FIFO queue at the DU. Bits that wait longer than $D_i$ slots are dropped and counted as \emph{deadline violations}. The budgets already respect the capacity, so the link itself never builds a queue. Congestion shows up instead as queueing and drops at the DU. Our primary metric was fixed before any tuning. It is the priority-weighted violation ratio
\begin{equation}
  V_w = \frac{\sum_i w_i\,\mathrm{violated}_i}{\sum_i w_i\,\mathrm{arrived}_i},
\end{equation}
in which a dropped LL bit counts three times. Alongside it we report per-class violations, link utilisation, unused allocation, DU delay, Jain fairness and decision time.

\section{Key Idea 1: the Deadline-Feasible Rate}
\label{sec:rstar}

What should a forecaster predict? The traffic volume is the obvious choice, but it is the wrong one, since a 2~ms cell hit by a burst needs far more than its average rate. We instead forecast the smallest constant rate at which no bit misses its deadline.

\begin{definition}[Deadline-feasible rate]
For arrivals $a_1,\dots,a_H$ and deadline $D$ slots,
\begin{equation}
  \rstar(a, D) = \max_{1\le s\le e\le H} \frac{\sum_{t=s}^{e} a_t}{e - s + 1 + D}.
  \label{eq:rstar}
\end{equation}
\end{definition}

Eq.~\eqref{eq:rstar} is the service-curve condition of network calculus~\cite{leboudec2001}. It says that every window of arrivals must be served within its own length plus the deadline. The condition is tight, and our tests confirm this to within one PRB-layer. When the deadline is long, bursts can be smoothed out and $\rstar$ approaches the mean rate. When it is short, $\rstar$ approaches the peak burst rate. Take a single burst of 8 units with an average of 1 per slot. It needs $\rstar = 8/2 = 4$ for $D=1$ slot, but only $8/8 = 1$ for $D=7$ slots (Figure~\ref{fig:rstar}). In our 2~ms cells $\rstar$ is usually 2--3$\times$ the mean, and its distribution is strongly right-skewed.

\begin{figure}[H]
\centering
\begin{tikzpicture}[font=\footnotesize,>=Stealth,xscale=0.75,yscale=0.42]
  \draw[->] (0,0) -- (13.5,0) node[right]{slot};
  \draw[->] (0,0) -- (0,9.2) node[above]{bits};
  \foreach \x/\h in {1/0,2/0,3/8,4/0,5/0,6/0,7/0,8/0} \draw[fill=iitblue!40] (\x-0.4,0) rectangle (\x+0.4,\h);
  \node[iitblue] at (3,8.6) {burst $=8$};
  \draw[very thick,accent] (3-0.5,4) -- (4.5,4) node[right,font=\scriptsize]{$D=1$: $\rstar = 8/2 = 4$};
  \draw[very thick,good!60!black] (3-0.5,1) -- (10.5,1) node[right,font=\scriptsize]{$D=7$: $\rstar = 8/8 = 1$};
\end{tikzpicture}
\caption{The deadline-feasible rate of one burst for a short and a long deadline: a short deadline forces a rate close to the peak.}
\label{fig:rstar}
\end{figure}

All deadline-aware controllers request the same quantity,
\begin{equation}
  \text{need}_i = B_i + \rstar_i, \qquad B_i = \frac{\text{backlog}_i}{\text{slack}_i},
\end{equation}
which is the rate needed to clear the current queue before its deadline, plus $\rstar_i$ for traffic that has not arrived yet. The controllers differ only in how they estimate $\rstar_i$ and in how they split a shortage. A single number is a poor basis for reserving a skewed quantity, so we forecast a \emph{range}.

\section{Key Idea 2: a Quantile Forecaster}
\label{sec:gbm}

The target is $y_i = \rstar(d_i[t-\tau : t+T], D_i)/\mathrm{cap}_i$, which is the deadline-feasible rate of the coming horizon divided by the cell's peak. All 28 features are computed from past slots. They are the last six horizon means (with their mean, standard deviation and maximum), two exponentially weighted averages, the last four slot demands, $\rstar$ of the last three horizons, the relative spectral efficiency, $D_i/H$ and a one-hot cell identifier. Normalisation uses configuration constants only. For each level $\alpha \in \{0.05, 0.25, 0.5, 0.75, 0.9, 0.95\}$ we train one histogram gradient-boosting quantile regressor (300 trees, 15 leaves) on 101k rows from 16 training traces. Training takes 12~s on a CPU. The training, validation and test seeds do not overlap (1000--1015, 2000--2003 and 3000--3019).

\section{Key Idea 3: the KKT (Constrained Newsvendor) Rule}
\label{sec:kkt}

\subsection{Objective and solution}

Suppose each cell comes with a forecast distribution $F_i$ of $\rstar_i$. We then maximise the expected priority-weighted need served in the interval,
\begin{equation}
  \max_{r}\ \sum_i w_i\,\mathbb{E}\big[\min(B_i + \rstar_i,\ r_i)\big]
  \quad \text{s.t.}\ \sum_i r_i \le C,\ \ 0\le r_i\le \mathrm{cap}_i .
  \label{eq:nv}
\end{equation}
The function $\mathbb{E}[\min(B_i + \rstar_i, r)]$ is concave in $r$. Its derivative is $w_i\,(1 - F_i(r - B_i))$, the weighted probability that the last bit of budget is actually used. Because the problem is separable and concave, the KKT conditions give its \emph{global} optimum in closed form:
\begin{equation}
  \boxed{\; r_i(\lambda) = \mathrm{clip}\Big(B_i + F_i^{-1}\big(1 - \lambda/w_i\big),\ 0,\ \mathrm{cap}_i\Big)\;}
  \label{eq:kkt}
\end{equation}
Only one scalar has to be found. It is $\lambda \ge 0$, the shadow price of link capacity, and a 60-step bisection chooses it so that $\sum_i r_i(\lambda) = C$. If the link is not full, $\lambda = 0$ and every cell gets its full predictive support. We interpolate $F_i^{-1}$ piecewise-linearly through the six forecast quantiles. The cost is $O(NK + N n_{\rm iter})$, which comes to a few tens of microseconds for eight cells.

\paragraph{Interpretation.} Eq.~\eqref{eq:kkt} cuts every cell at the \emph{same priority-weighted risk}. A cell with a wide forecast or a high priority gets a larger margin above its backlog, while a cell with a confident forecast gets roughly its median. Nobody tunes these margins by hand. They are \emph{derived} from the forecast and the link constraint.

\paragraph{Toy example.} Say the bisection returns $\lambda = 0.6$. An LL cell ($w=3$) is then funded up to the $1 - 0.6/3 = 0.8$ quantile of its forecast, and an eMBB cell ($w=1$) up to its $1-0.6 = 0.4$ quantile (Figure~\ref{fig:kkt-toy}).

\begin{figure}[H]
\centering
\begin{tikzpicture}[font=\footnotesize,>=Stealth,yscale=0.9]
  \draw[->] (0,0) -- (9.6,0);
  \draw[->] (0,0) -- (0,4.6) node[above]{budget $r_i$};
  % cells
  \foreach \x/\b/\m/\lab/\col in {1/0.8/2.4/LL 1/accent,2.6/0.5/2.0/LL 2/accent,4.2/1.3/0.6/eMBB 1/iitblue,5.8/0.4/0.9/eMBB 2/iitblue,7.4/0.9/0.4/eMBB 3/iitblue}{
    \draw[fill=gray!35] (\x-0.45,0) rectangle (\x+0.45,\b);
    \draw[fill=\col!30] (\x-0.45,\b) rectangle (\x+0.45,\b+\m);
    \node[below] at (\x,0) {\lab};
  }
  \node[align=left,anchor=west,font=\scriptsize] at (8.1,3.9) {\tikz\draw[fill=gray!35] (0,0) rectangle (0.3,0.2); backlog $B_i$};
  \node[align=left,anchor=west,font=\scriptsize] at (8.1,3.4) {\tikz\draw[fill=accent!30] (0,0) rectangle (0.3,0.2); LL: $F^{-1}(0.8)$};
  \node[align=left,anchor=west,font=\scriptsize] at (8.1,2.9) {\tikz\draw[fill=iitblue!30] (0,0) rectangle (0.3,0.2); eMBB: $F^{-1}(0.4)$};
  \node[align=left,anchor=west,font=\scriptsize] at (8.1,2.2) {$\textstyle\sum_i r_i = C$\\one price $\lambda = 0.6$};
\end{tikzpicture}
\caption{Illustration of the KKT rule: every cell first receives its backlog rate, then a quantile of its forecast whose level is set by its priority and a common capacity price $\lambda$.}
\label{fig:kkt-toy}
\end{figure}

\subsection{Properties}

We prove four properties of the rule. The proofs are in Appendix~\ref{app:proofs}.

\begin{proposition}[Feasibility]\label{prop:feas}
For any forecast and any $\lambda$, the budgets satisfy the caps, and the bisection output satisfies $\sum_i r_i \le C$: the link is never overloaded.
\end{proposition}

\begin{proposition}[Equal weighted tail probability]\label{prop:equal}
At an interior optimum, $w_i\big(1 - F_i(r_i - B_i)\big) = \lambda$ for every cell. With point forecasts (degenerate $F_i$), the rule reduces to a priority-ordered fill.
\end{proposition}

\begin{proposition}[Backlog first]\label{prop:backlog}
Any cell with $\lambda < w_i$ receives at least $\min(B_i,\mathrm{cap}_i)$ before the uncertain demand of any cell of equal or lower weight is funded.
\end{proposition}

\begin{proposition}[Robustness to miscalibration]\label{prop:robust}
If $\sup_x |F_i(x) - G_i(x)| \le \epsilon$ for every cell, where $G_i$ is the true distribution, then
\begin{equation}
  J_G(r^G) - J_G(r^F) \le \epsilon\sum_i w_i\big[(r_i^G - B_i)^+ + (r_i^F - B_i)^+\big] \le 2\,\epsilon\, w_{\max}\, C .
\end{equation}
The loss due to a miscalibrated forecast grows at most linearly in $\epsilon$.
\end{proposition}

\subsection{Control loop}

Figure~\ref{fig:loop} shows the full loop, which runs every 10~ms. We handle the edge cases explicitly. If a forecast is not finite, the controller falls back to window quantiles. A cell with zero demand gets zero budget plus a share of any headroom. When the link is not binding, every cell gets its full predictive support.

\begin{figure}[H]
\centering
\begin{tikzpicture}[font=\footnotesize,>=Stealth,node distance=0.5cm and 0.45cm,
   st/.style={draw,rounded corners,fill=iitblue!10,align=center,minimum height=1.1cm,text width=2.55cm}]
  \node[st] (obs) {Delayed state ($t-\tau$):\\queues, HOL age,\\history, $\eta$};
  \node[st,right=of obs] (feat) {28 features\\per cell};
  \node[st,right=of feat] (gbm) {Forecaster:\\quantiles of $\rstar_i$};
  \node[st,right=of gbm,fill=accent!12] (kkt) {Allocation rule\\(KKT / BPN)};
  \node[st,below=1.1cm of kkt,fill=good!12] (du) {DU schedulers:\\$\lfloor r_i/\beta\rfloor$ PRB-layers\\per slot};
  \node[st,left=of du] (link) {Shared link\\(never overloaded)};
  \node[st,left=of link] (q) {FIFO queues,\\drops after $D_i$};
  \node[st,left=of q] (traffic) {Traffic\\$a_i(t)$, $\eta_i(t)$};
  \draw[->,thick] (obs) -- (feat);
  \draw[->,thick] (feat) -- (gbm);
  \draw[->,thick] (gbm) -- (kkt);
  \draw[->,thick] (kkt) -- node[right,font=\scriptsize]{$r_i$ held $T$ slots} (du);
  \draw[->,thick] (du) -- (link);
  \draw[->,thick] (link) -- (q);
  \draw[->,thick] (traffic) -- (q);
  \draw[->,thick,dashed] (q.north) to[out=110,in=-70] (obs.south);
  \node[above=0.08cm of kkt,text=iitblue] {every $T = 10$\,ms};
\end{tikzpicture}
\caption{The closed control loop of Part~B. The forecaster and the allocation rule are the two blocks developed in this chapter.}
\label{fig:loop}
\end{figure}

\section{Extended Study: Theory and the Strongest Baselines}
\label{sec:ext}

\subsection{The control problem as a POMDP}

Let the state be $s(t) = (Q(t), Z(t), \eta(t))$, which collects the queue age profiles, the hidden traffic state and the spectral efficiencies. At epochs $t_k = kT$ the controller solves
\begin{equation}
  \min_{\pi}\ \mathbb{E}_\pi\Big[\sum_k \sum_{t=t_k}^{t_k+T-1} \sum_i w_i\,\mathrm{dropped}_i(t)\Big],
  \qquad r(t_k) = \pi\big(s(t_k - \tau),\text{history}\big).
\end{equation}
This partially observable MDP cannot be solved directly. Its continuous state has hundreds of dimensions, the traffic law is unknown, and each action is held for $T$ slots. We implemented the three main families of approximations instead:
\begin{enumerate}
  \item \textbf{One-step stochastic programs.} The KKT newsvendor rule (Section~\ref{sec:kkt}) applied to a learned or model-based distribution, including a Markov chain fitted online.
  \item \textbf{Lyapunov drift.} MaxWeight~\cite{tassiulas1992} and \textbf{predictive MaxWeight}~\cite{huang2016}. The latter serves cells in order of $w_i\,(Q_i + \ell H m_i)$, the queue the cell \emph{will} have given forecast median $m_i$, and gives each cell up to $Q_i/T + (1+\mu)\,m_i$.
  \item \textbf{Direct policy search.} A class-level policy in which each cell reserves $B_i + F_i^{-1}(p_{c(i)})$ and any shortage is shared by weighted water-filling with a learned LL weight. Its three parameters $(p_{\rm LL}, p_{\rm eMBB}, w_{\rm LL})$ are learned on training seeds with the cross-entropy method (CEM)~\cite{rubinstein1999}. One can think of it as a small and reliable analogue of a deep-RL xApp.
\end{enumerate}

\subsection{Markov-fluid theory of the deadline-feasible rate}

Consider Markov-modulated traffic with phase chain $P$ (stationary distribution $\pi$), per-phase log-MGF $\psi_k(\theta)$ and $\Phi(\theta) = \mathrm{diag}(e^{\psi_k(\theta)})$. For this model we derived the following results.
\begin{itemize}
  \item the window log-MGF $\Lambda_l(\theta) = \log\big(\pi^{\mathsf T}\Phi(P\Phi)^{l-1}\mathbf 1\big)$ and the effective bandwidth $\alpha(\theta) = \log\rho(P\Phi(\theta))/\theta$~\cite{kesidis1993,kelly1996};
  \item a \textbf{Chernoff--union bound} on the deadline-feasible rate,
  \begin{equation}
    \Pr(\rstar > x) \le \sum_{l=1}^{H}(H-l+1)\,\min_\theta \exp\big(-\theta x(l+D) + \Lambda_l(\theta)\big);
  \end{equation}
  \item a \textbf{large-deadline decay} $\limsup_{D\to\infty} D^{-1}\log\Pr(\rstar > x) \le -x\,\theta^\ast(x)$, which links $\rstar$ to classical buffer-overflow asymptotics~\cite{anick1982,duffield1994};
  \item a pathwise \textbf{shortfall sandwich} valid for any arrivals,
  \begin{equation}
    (D+1)(\rstar - r)^+ \le V(r) \le \textstyle\sum_t (a_t - r)^+,
  \end{equation}
  where $V(r)$ is the volume violated at constant rate $r$. It follows that the newsvendor surrogate~\eqref{eq:nv} is a lower bound on the violated volume.
\end{itemize}
The bounds held in every validation run. They are, however, conservative by a factor of 25--100 at 2--10~ms deadlines. That gap is a strong reason to \emph{learn} the distribution of $\rstar$ from data instead of computing it from a model.

\section{B2: LCQT, a Bayesian Link-Aware Quantile Transformer}
\label{sec:lcqt}

\subsection{Architecture}

LCQT (Figure~\ref{fig:lcqt}) can replace the GBM forecaster directly. It is assembled from established components: PatchTST patches~\cite{nie2023}, iTransformer-style attention across variates~\cite{liu2024}, TFT-style multi-quantile output~\cite{lim2021}, FiLM conditioning~\cite{perez2018film}, deep ensembles~\cite{lakshminarayanan2017} and MC dropout~\cite{gal2016}. We made six changes to adapt it to the shared fronthaul link.
\begin{enumerate}
  \item \textbf{Deadline-aware patch tokens.} The 352 slots of history are cut into 16 patches of $H=22$ slots. Each patch token carries its own $\rstar$ at the cell's deadline, along with its mean and maximum. A separate summary token holds the 28 GBM features.
  \item \textbf{Deadline/priority FiLM.} A single network serves both LL and eMBB cells and is modulated by $(D_i/H, p_i)$.
  \item \textbf{Link-aware cross-cell attention.} After the temporal encoder, attention runs over the eight cells that share the link.
  \item \textbf{Non-crossing quantile head.} Writing $q_k = q_{k-1} + \mathrm{softplus}(a_k)$ makes the quantiles monotone by construction.
  \item \textbf{Bayesian mixture layer.} Three trained seeds $\times$ eight dropout passes form a posterior, and its \emph{mixture} distribution is what the allocation rule receives.
  \item \textbf{Differentiable KKT layer.} This lets the network be trained through the coupled capacity constraint.
\end{enumerate}
The model has 63k parameters (44k without the cross-cell layer). It uses the same seeds and target as the GBM and trains in about 6 minutes on a 4-core CPU, with 40 epochs of AdamW on a one-cycle schedule and early stopping on validation pinball loss.

\begin{figure}[H]
\centering
\begin{tikzpicture}[font=\scriptsize,>=Stealth,
    box/.style={draw=iitblue,fill=soft,rounded corners=2pt,align=center,minimum height=1.0cm,text width=1.75cm,inner sep=2pt},
    hot/.style={box,draw=accent,fill=accent!8},
    grn/.style={box,draw=good,fill=good!8}]
  \node[box] (hist) {Per-cell history\\$a_i(t)$, 352 slots};
  \node[box,right=0.3cm of hist] (patch) {16 patches\\+ $\rstar$, mean, max};
  \node[box,below=0.35cm of patch] (feat) {Summary token\\28 features};
  \node[box,right=0.3cm of patch] (enc) {Temporal encoder\\2 layers, $d{=}48$\\4 heads};
  \node[box,right=0.3cm of enc] (film) {FiLM\\$(D_i/H,\ p_i)$};
  \node[hot,right=0.3cm of film] (cross) {Cross-cell\\attention\\(8 cells)};
  \node[box,right=0.3cm of cross] (head) {Non-crossing\\6 quantiles};
  \node[grn,below=1.6cm of cross,text width=2.3cm] (bayes) {Bayesian layer\\3 seeds $\times$ 8 dropout\\$\Rightarrow$ mixture $\hat F_i$};
  \node[grn,left=0.6cm of bayes,text width=2.9cm] (kkt) {KKT layer\\$r_i = B_i + \hat F_i^{-1}(1-\lambda^\ast/w_i)$};
  \node[hot,left=0.6cm of kkt,text width=3.2cm] (loss) {Decision loss\\$\sum_i p_i (D_i{+}1)(B_i{+}\rstar_i - r_i)^+$\\$+\ \beta\,$pinball};
  \draw[->] (hist) -- (patch);
  \draw[->] (patch) -- (enc);
  \draw[->] (feat.east) -| ([xshift=-0.2cm]enc.south);
  \draw[->] (enc) -- (film);
  \draw[->] (film) -- (cross);
  \draw[->] (cross) -- (head);
  \draw[->] (head.south) |- (bayes.east);
  \draw[->] (bayes) -- (kkt);
  \draw[->,dashed,accent] (kkt) -- node[above]{$\partial\lambda^\ast$} (loss);
\end{tikzpicture}
\caption{LCQT: link-aware cross-cell quantile Transformer with the Bayesian mixture layer and the differentiable KKT layer used for decision-focused training.}
\label{fig:lcqt}
\end{figure}

\subsection{Bayesian mixture layer}

The posterior samples $s = 1,\dots,MS$ (ensemble members $\times$ dropout passes) each give a quantile function $Q_i^{(s)}$. Averaging their quantiles would throw away the disagreement between samples. We pass the KKT rule the mixture distribution instead,
\begin{equation}
  \hat F_i(x) = \frac{1}{MS}\sum_s F_i^{(s)}(x).
\end{equation}
so a cell whose forecast the posterior is unsure about gets a wider margin at the same price $\lambda^\ast$. Dropout is seeded per decision slot, which keeps runs exactly reproducible.

\subsection{Differentiable KKT layer and decision-focused training}

The forward pass finds $\lambda^\ast$ by bisection. For the backward pass we apply the implicit function theorem to the binding link constraint $g(\lambda,q) = \sum_i r_i(\lambda,q) - C = 0$, which gives
\begin{equation}
  \frac{\partial\lambda^\ast}{\partial q} = -\frac{\partial g/\partial q}{\partial g/\partial \lambda},
\end{equation}
and in code this is one Newton step from the detached root. The decision-focused model (DF-LCQT) is trained on the shortfall-sandwich lower bound of the dropped bits, $\sum_i p_i (D_i+1)(B_i + \rstar_i - r_i)^+$, with a pinball regulariser added. The contexts $(B_i, w_i)$ depend on the policy. For that reason we log them in closed loop on training seeds and refresh them with one DAgger round. Unit tests check that the layer agrees with the numpy allocator to $10^{-6}$. They also check that $\partial \sum_i r_i / \partial q = 0$ when the link is full, since budget can then only move between cells.

\section{B3: The Back-Pressure Newsvendor}
\label{sec:bpn}

\subsection{Motivation}

A budget acts as a per-cell cap. If capacity is reserved for the upper quantile of a 2~ms cell and then goes unused, no other cell can have it. A pure reservation rule protects LL cells well for this reason, but it leaves the eMBB cells short. Back-pressure (MaxWeight) fills the link in order of queue pressure using median demands. It wastes less capacity, yet it ignores both deadlines and forecast uncertainty. The back-pressure newsvendor (BPN) puts the two ideas together. It \emph{fills a base in back-pressure order and then spends the leftover capacity on newsvendor margins} (Figure~\ref{fig:bpn}).

\begin{figure}[H]
\centering
\begin{tikzpicture}[font=\footnotesize,>=Stealth,
   s/.style={draw,rounded corners,align=left,text width=6.2cm,inner sep=5pt}]
  \node[s,fill=iitblue!8] (st1) {\textbf{Stage 1: base in back-pressure order}\\[2pt]
    base rate $b_i = a\,Q_i/T + (1-a)\,B_i$\\
    base demand $d_i = b_i + (1+\mu)\,\hat F_i^{-1}(p^0_{c(i)})$\\
    shortage shared $\propto v_{c}\,(\pi_i/\bar\pi)^{\eta}$\\
    pressure $\pi_i = Q_i + \ell H m_i$\\[2pt]
    \scriptsize $\eta = 0$: class-weighted water-fill (reservation)\\ $\eta = \eta_{\max}$: strict back-pressure order};
  \node[s,fill=good!10,right=0.8cm of st1] (st2) {\textbf{Stage 2: leftover by the KKT rule}\\[2pt]
    $r_i = \max\!\big(r_i^{(1)},\ b_i + \hat F_i^{-1}(1 - \lambda/\omega_{c(i)})\big)$\\
    with $\lambda$ such that $\sum_i r_i = C$\\[2pt]
    \scriptsize newsvendor margins only where the forecast is wide};
  \draw[->,very thick] (st1) -- node[above,font=\scriptsize]{$r^{(1)}$} (st2);
  \node[draw,rounded corners,fill=orange!10,below=0.6cm of st1.south east,anchor=north,xshift=0.4cm,text width=12.8cm,align=center,inner sep=4pt] (cem) {\textbf{Learning:} 8 parameters $\theta = (a, \mu, p^0_{\rm LL}, p^0_{\rm eMBB}, v_{\rm LL}, \eta, \ell, \omega_{\rm LL})$ by CEM on \emph{training} seeds, started from the tuned optimum of each contained baseline};
\end{tikzpicture}
\caption{The two stages of the back-pressure newsvendor and how its parameters are learned.}
\label{fig:bpn}
\end{figure}

\subsection{The policy}

\paragraph{Stage 1.} The base rate $b_i = a\,Q_i/T + (1-a)\,B_i$ blends MaxWeight's queue drain rate with the deadline-aware backlog rate. To get each cell's base demand we add a class-specific base quantile of its forecast, $d_i = b_i + (1+\mu)\hat F_i^{-1}(p^0_{c(i)})$. If the base demands exceed the link, the shortage is shared in proportion to $v_{c}(\pi_i/\bar\pi)^{\eta}$. Here the pressure is $\pi_i = Q_i + \ell H m_i$, the queue plus look-ahead $\times$ forecast median. Changing the exponent $\eta$ moves the policy smoothly from class-weighted water-filling ($\eta = 0$, reservation) to strict back-pressure ordering ($\eta = \eta_{\max}$).

\paragraph{Stage 2.} The KKT rule distributes whatever capacity is left after stage~1. Each cell gets $r_i = \max(r_i^{(1)},\ b_i + \hat F_i^{-1}(1 - \lambda/\omega_{c(i)}))$, and bisection sets $\lambda$ so that $\sum_i r_i = C$.

\subsection{Both strongest baselines are special cases}

For suitable parameters BPN \emph{is} each of the strongest baselines. Unit tests verify this exactly under shortage. The settings are:
\begin{itemize}
  \item predictive MaxWeight: $a = 1$, $p^0 = 0.5$, $\mu$ = its margin, $\eta = \eta_{\max}$;
  \item the CEM class-level policy: $a = 0$, $p^0$ = its class levels, $v_{\rm LL} = w_{\rm LL}$, $\eta = 0$.
\end{itemize}
A search over the BPN parameters that starts from the tuned optimum of each baseline can then only match or improve on both of them.

\subsection{Learning with equal search effort}

We learn the eight parameters with CEM on training seeds only, using 8 episodes per candidate, which is the same search used for the baselines. The search is started twice, once at the tuned optimum of each contained baseline. For a fair comparison, both baselines were also \emph{re-tuned with the same CEM search} (``CEM-tuned''), so every controller had equal tuning effort. Two versions are reported. BPN~v1 has 6 parameters and always uses strict back-pressure in stage~1. BPN~v2 has 8 parameters, adding the sharing exponent $\eta$ and the margin $\mu$, and it contains the reservation policy as well.

\section{Summary}

This chapter developed the Part~B controllers. The deadline-feasible rate turns deadlines into a demand target, and a calibrated quantile forecaster describes how uncertain that target is. The closed-form KKT rule then splits the link at equal priority-weighted risk, and we proved four properties for it. In the extended study we framed the problem as a POMDP, derived Markov-fluid bounds and introduced the strongest baselines. LCQT contributes a link-aware Bayesian Transformer forecaster with a differentiable KKT layer. The back-pressure newsvendor brings reservation and back-pressure into one learned policy class that contains both strongest baselines. Chapter~\ref{ch:partb-results} evaluates all of these methods.
